\documentclass[11pt,a4paper]{article}
\usepackage[margin=25mm]{geometry}
\usepackage{amsmath,amssymb,amsthm,booktabs}
\newtheorem{theorem}{Theorem}
\title{The exceptional Coxeter censuses}
\author{Human Compression project: F006 recovery}
\date{Working manuscript; distinct review pending}
\begin{document}\maketitle
We allow an adjacent product of order two in a string C-marking.
Consequently its string diagram need not be connected. Counts below
identify a marking with its reverse; parenthesized numbers count
self-dual classes. The complete proofs of enumeration by stabilizers
of initial strings and of its implementation through double cosets
are acquired prerequisites. The exact finite executions specified
below are additional premises. The classical finite Coxeter
presentations and orders, and exact GAP group operations, are retained
inputs throughout.

\begin{theorem}\label{census}
The complete census for the following five finite Coxeter groups is
given in the table. Every omitted rank has zero classes.
\[
\begin{array}{c|rrrrr|r}
 &3&4&5&6&7&\text{total}\\\hline
 F_4&3(0)&2(1)&0&0&0&5(1)\\
 H_4&45(2)&14(4)&0&0&0&59(6)\\
 E_6&87(12)&50(4)&10(2)&0&0&147(18)\\
 E_7&1577(10)&1987(0)&741(0)&167(0)&18(0)&4490(10)\\
 E_8&6746(117)&3584(22)&986(2)&310(0)&63(1)&11689(142)
\end{array}
\]
All markings for $F_4,H_4,E_6,E_8$ are connected. The connected
part of the $E_7$ census is
\[
\begin{array}{c|rrrrr|r}
 &3&4&5&6&7&\text{total}\\\hline
 E_7&1577(10)&1525(0)&465(0)&95(0)&0&3662(10).
\end{array}
\]
Thus the respective maximum ranks are $4,4,5,7,7$ under the
stated convention; the connected maximum for $E_7$ is $6$.
\end{theorem}
\begin{proof}
The models are exact reflection permutations on finite root sets.
The simple reflections satisfy the appropriate Coxeter relations,
and their generated groups have orders
\[
 1152,\quad14400,\quad51840,\quad2903040,\quad696729600.
\]
The classical presentations therefore give surjections from the
specified Coxeter groups onto the models. Equality with the
classical orders makes these maps isomorphisms. The complete
constructors, root coordinates, permutations and identity checks
are retained finite evidence.

For $F_4,H_4,E_6$, the enumeration uses the full automorphism
action on the complete involution set. For $E_7$, the exact guard
\[
 |E_7:Z(E_7)|=|\operatorname{Aut}(E_7)|
\]
allows conjugation by the group itself. Its centralizers are the
inverse images of the pointwise automorphism stabilizers. The
central kernel acts trivially and changes none of their orbits.

For $E_8$, use the acquired construction on two permutation
domains. In the original degree-$240$ model, let $z$ be the
nonidentity central element and let $\chi$ be permutation parity.
The finite guards establish $|z|=2$ and $\chi(z)=0$.
Thus $\phi(g)=z^{\chi(g)}g$ is an involutory automorphism.
Represent $g$ on two copies of the original domain by $g$ and
$\phi(g)$, respectively. The permutation interchanging the copies
implements $\phi$. Adjoining it gives a group $A$ normalizing
the degree-$480$ copy of $E_8$. The exact guard gives
\[
 |A|=1393459200,\quad |C_A(E_8)|=2,
 \quad |A:C_A(E_8)|=|\operatorname{Aut}(E_8)|=696729600.
\]
Hence conjugation by $A$ supplies the full automorphism group.
The acquired proof justifies its use at every prefix stabilizer
and in every reversal test; a full action is checked here, rather
than presumed from the form of the central automorphism.

The $E_8$ execution enumerates candidate orbits through classes
in the current centralizer and the corresponding double cosets.
This is the acquired exact implementation of the same stabilizer
search. It also discards central involutions in the derived
subgroup. To justify this last step, if a distinguished central
involution $t_i$ occurred in a C-marking, the intersection property
would give
\[
 G=\langle t_i\rangle\times
       \langle t_j:j\ne i\rangle.
\]
Every commutator lies in the second factor, whereas $t_i$ does not.
Thus such a mark lies outside $G'$. The subgroup $Z(G)\cap G'$
is characteristic, so deleting its involutions preserves the
candidate-domain invariance and excludes no valid marking.
The acquired completion test uses the subgroup generated by all
remaining eligible elements, as calculated by the normal closure
of their complete class representatives.

The complete ordered counts at successive nonempty ranks are
$6,3$ for $F_4$; $88,24$ for $H_4$; $162,96,18$ for $E_6$;
$3144,3974,1482,334,36$ for $E_7$; and
$13375,7146,1970,620,125$ for $E_8$.
The corresponding self-dual counts are
$(0,1)$, $(2,4)$, $(12,4,2)$, $(10,0,0,0,0)$ and
$(117,22,2,0,1)$.
The exact test for an automorphism sending a tuple to its reverse
gives these fixed-point counts. Reversal is an involution on the
ordered automorphism classes, so $N$ ordered classes and $S$
fixed classes give $(N+S)/2$ classes modulo duality. This proves
the first table.

Restricting to tuples whose adjacent products have order greater
than two leaves every tuple except the disconnected $E_7$
markings. In that group it leaves $3144,3050,930,190,0$ ordered
tuples at ranks three through seven, with self-dual counts
$10,0,0,0,0$. The same reversal formula proves the second table.
The five complete executions cover every rank through
$\lfloor\log_2|G|\rfloor$, namely $10,13,15,21,29$.
They retain the actual generating tuples, action and model
generators, all rank summaries, and the final exhaustion records.
The acquired enumeration proof establishes coverage and
uniqueness, including every zero row. The asserted maxima follow
from the last nonzero ranks.
\end{proof}

\paragraph{A separate rank-seven witness.}
The first retained rank-seven $E_7$ tuple has adjacent product
orders $(6,6,4,3,3,2)$. Its final mark is central of order two.
The other six marks generate a subgroup of order $1451520$
intersecting that central subgroup trivially; together they
generate the reflection group of order $2903040$. A separate
execution verifies the intersection property directly for all
$16384$ pairs of subsets of the seven marks. The literal tuple
and execution are retained. This provides a witness independently
of the census counts. Its diagram is disconnected because the
last adjacent product has order two.

\paragraph{Comparison with the source.}
Chen and Rowley, \emph{Abstract Regular Polytopes of Finite
Irreducible Coxeter Groups}, Electronic Journal of Combinatorics
32(4) (2025), P4.18, Table~1, physical page~12, gives exactly
the connected counts in Theorem~\ref{census}. The displayed
definition on page~1 permits disconnected strings, and page~2
expressly discusses them for $B_n$. Under that literal convention,
the $E_7$ clause of Theorem~3(ii) requires the rank-seven
correction. Its rank-six assertion is correct with connectedness
imposed. The independently acquired $I_2(m)$ and $H_3$ component
records the analogous scope corrections for those groups.

The original authors supply the target results and published
connected counts. The enumeration method follows the established
Hartley--Hulpke and Leemans--Mixer approach, with attribution and
the precise executable tests given in the acquired proofs. The
two-domain implementation and central exclusion use standard
group constructions; no historical priority is claimed. The new
executions recover the printed counts and retain disconnected
cases under the source's literal definition. They use GAP as an
explicit software premise and are not Lean certificates.
The present manuscript makes no open-conjecture or numerical
compression claim. The acquired uniform $D_n$ proofs and the
small-group component are required for a complete F006 recovery.
\end{document}
